\begin{figure}[t!]
	\begin{center}
		\begin{tikzpicture}[
			framed,
			every node/.style={scale=0.8},
			modal,%
			node distance =.6cm and .8cm%
			]
			\node[world] (n1)
			[label=150:{\footnotesize$n$}]
			{$pqrs$};
			
			\node[world] (n2)
			[label=180:{\footnotesize$m$}, below left=of n1]
			{$pqrs$};
			
			\node[world] (n3)
			[below right=of n1]
			{$r$};
			
			\node[world] (n4)
			[below=of n2]
			{$r$};
			
			\node[world] (n5)
			[left=of n4]
			{$r$};
			
			\node[world] (n6)
			[right=of n4]
			{$r$};
			
			\node[world] (n7)
			[fill=gray!15, label=150:{\footnotesize$n$}, right=3cm of n1]
			{$pqr$};
			
			\node[world] (n8)
			[fill=gray!15, label=180:{\footnotesize$m$}, below= of n7]
			{$r$};
			
			\node (a)
			[scale=1.2, below right=1.3cm of n4]
			{\hspace*{-1.3cm}(a)};
			\node
			[scale=1.2, below=.2cm of n8]
			{(b)};
			
			\path[->]
			(n1)
			edge[bend left]
			node[midway, right]{\footnotesize$\mathsf{b}$}
			(n2);
			
			\path[->]
			(n1)
			edge[bend right]
			node[midway, left]{\footnotesize$\mathsf{a}$}
			(n2);
			
			\path[->]
			(n1)
			edge
			node[midway, right]{\footnotesize$\mathsf{c}$}
			(n3);
			
			\path[->]
			(n2)
			edge
			node[midway, right]{\footnotesize$\mathsf{b}$}
			(n4);
			\path[->]
			(n2)
			edge
			node[midway, right]{\footnotesize$\mathsf{c}$}
			(n6);
			\path[->]
			(n2)
			edge
			node[midway, left]{\footnotesize$\mathsf{a}$}
			(n5);
			
			% \path[-]
			% (n3)
			% edge[dotted,bend left]
			% node[midway, right]{\footnotesize$\neq_\cmp$}
			% (n6);
			\path[-]
			(n4)
			edge[dotted]
			node[midway, above]{\footnotesize$=_\cmp$}
			(n5);
			\path[-]
			(n2)
			edge[dotted]
			node[midway, below]{\footnotesize$=_\cmp$}
			(n3);
			\path[-]
			(n5)
			edge[dotted,bend right=40]
			node[midway, below]{\footnotesize$=_\cmp$}
			(n6);
			
			\path[->]
			(n7)
			edge[bend left]
			node[midway, right]{\footnotesize$\mathsf{b}$}
			(n8);
			\path[->]
			(n7)
			edge[bend right]
			node[midway, left]{\footnotesize$\mathsf{a}$}
			(n8);
		\end{tikzpicture}
	\end{center}
	\caption{Model and Selected Model}\label{fig:model:selected}
	\medskip\medskip\medskip
\end{figure}

\begin{example}\label[example]{ex:selection-tbox}
 	   Consider the TBox
        \[\tbox=\set{
            p\doteq q\land r,
            q\doteq \tup{\apath =_\cmp \ppath{b}},
            r \doteq [\ppath{bc}\neq_\cmp\ppath{c}],
            s\doteq \tup{\apath =_\cmp \ppath{c}}
        },\]
    and let $\model$ be the model in \Cref{fig:model:selected}(a).
    In this figure, dotted arrows represent the relevant equality information (i.e., $\data_\cmp$ applied to those nodes returns equal values, and returns distinct values if applied to any other two different nodes).
%    For the sake of clarity, we have chosen to avoid drawing self-loop dotted arrows on the nodes of the model.
    Immediately, $\model \vDash \tbox$. 
    For instance, ${\truthset{q}{\model} = \{n,m\} = \truthset{\tup{\apath =_\cmp \ppath{b}}}{\model}}$.
    Similarly, if we take $\N$ to be the set of all nodes in $\model$, ${\truthset{r}{\model} = \N = \truthset{[\ppath{bc}\neq_\cmp\ppath{c}]}{\model}}$.
    Notice also that this model establishes that $p$ is satisfiable w.r.t.\ $\tbox$.
    This is easily verified since $n \in \truthset{p}{\model} = \truthset{q \land r}{\model} = \truthset{q}{\model} \cap \truthset{r}{\model} = \{n,m\}$.

    In turn, the model in \Cref{fig:model:selected}(b) corresponds to $\model{\restriction}_{\sel(p,n)}$.
	Intuitively, $\model{\restriction}_{\sel(p,n)}$ can be understood as the ``relevant part'' of $\model$ which is necessary to check the satisfiability of~$p$.
	It is worth pointing out that such a  ``relevant part'' is not simply a restriction, it carries with it a change in valuations.
	For instance, ${\truthset{s}{\model} = \truthset{\tup{\apath =_\cmp \ppath{c}}}{\model}} = \{n,m\}$, whereas ${\truthset{s}{\model{\restriction}_{\sel(p,n)}} = \truthset{\tup{\apath =_\cmp \ppath{c}}}{\model{\restriction}_{\sel(p,n)}}} = \emptyset$.
	Notice however, that this change in valuations does not affect the satisfiability of $p$ w.r.t.~$\tbox$.
	More precisely, it can be checked that $n \in \truthset{p}{\model{\restriction}_{\sel(p,n)}}$ since $\truthset{p}{\model{\restriction}_{\sel(p,n)}} = \truthset{q \land r}{\model{\restriction}_{\sel(p,n)}} = \truthset{q}{\model{\restriction}_{\sel(p,n)}} \cap \truthset{r}{\model{\restriction}_{\sel(p,n)}} = \{n\}$.
\end{example}
